\documentclass[12pt,a4paper,reqno]{amsart}

% language
\usepackage[greek,english]{babel}
\usepackage[utf8]{inputenc}
\usepackage{alphabeta}

% change default names to greek
\addto\captionsenglish{
    \renewcommand{\contentsname}{Περιεχόμενα}
    \renewcommand{\refname}{Βιβλιογραφία}
    \renewcommand{\datename}{Ημερομηνία:}
    \renewcommand{\urladdrname}{Ιστοσελίδα}
}

% math 
\usepackage{amsmath,amsthm,amssymb,amscd}

% font
\usepackage[cal=euler]{mathalfa}
\usepackage{libertinus-type1}
% \usepackage{txfonts} % for upright greek letters
\usepackage{bm} % for bold symbols
\usepackage{bbm} % for the simply-looking bb symbols

% miscellaneous 
\usepackage{changepage} %for indenting environments
\usepackage{csquotes} % example: \textcquote{}
\usepackage{blkarray}
\setcounter{MaxMatrixCols}{20} % default for pmatrix is 10!!
\usepackage{ytableau}
\usepackage{array} %needed to increase the vertical length in a tabular

% drawing
\usepackage{tikz,tikz-cd}
\usetikzlibrary{shapes.misc, patterns, matrix, calc, intersections,positioning,arrows.meta}
\usepackage{graphics,graphicx}
\usepackage{float} % provides enhanced control and customization options for floating objects such as figures and tables

% colors
\usepackage{xcolor}
\definecolor{darkcandyapplered}{rgb}{0.64, 0.0, 0.0}
\definecolor{midnightblue}{rgb}{0.1, 0.1, 0.44}
\definecolor{mylightblue}{HTML}{336699}
\definecolor{burntorange}{rgb}{0.8, 0.33, 0.0}
\definecolor{iceberg}{rgb}{0.44, 0.65, 0.82}
\definecolor{applegreen}{rgb}{0.55, 0.71, 0.0}
\definecolor{canaryyellow}{rgb}{1.0, 0.94, 0.0}
\definecolor{lightgray}{rgb}{0.83, 0.83, 0.83}

% hrefs
\usepackage{hyperref}
\usepackage[noabbrev,capitalize]{cleveref}
\hypersetup{
    pdftoolbar=true,        
    pdfmenubar=true,        
    pdffitwindow=false,     
    pdfstartview={FitH},  % fits the width of the page to the window
    pdftitle={},
    pdfauthor={},
    pdfsubject={},
    pdfkeywords={},
    pdfnewwindow=true,  % links in new window
    colorlinks=true,  % false: boxed links; true: colored links
    linkcolor=darkcandyapplered,   % color of internal links
    citecolor=midnightblue,  % color of links to bibliography
    urlcolor=cyan,  % color of external links
    linktocpage=true  % changes the links from the section body to the page number
    }

% geometry
\textwidth=16cm 
\textheight=21cm 
\hoffset=-55pt 
\footskip=25pt

% thm envs (you might need to change the path)
\theoremstyle{definition}
\newtheorem{theorem}{Θεώρημα}
\newtheorem{proposition}[theorem]{Πρόταση}
\newtheorem{lemma}[theorem]{Λήμμα}
\newtheorem{corollary}[theorem]{Πόρισμα}
\newtheorem{conjecture}[theorem]{Εικασία}
\newtheorem{problem}[theorem]{Πρόβλημα}
\newtheorem*{claim}{Ισχυρισμός}
\newtheorem{observation}[theorem]{Παρατήρηση}
\newtheorem{definition}[theorem]{Ορισμός}
\newtheorem{question}[theorem]{Ερώτημα}
\newtheorem*{questions}{Ερωτήματα}
\newtheorem*{que}{Ερώτημα}
\newtheorem{example}[theorem]{Παράδειγμα}
\newtheorem{exercise}{Άσκηση}
\newtheorem*{zorn}{Λήμμα του Zorn}
\newtheorem*{example_6.6}{Παράδειγμα~6.6}

\theoremstyle{remark}
\newtheorem*{remark}{Παρατήρηση}

% fixes the correct numbering of environments
\numberwithin{theorem}{section}
\numberwithin{exercise}{section}
\numberwithin{equation}{section}

% math ops 
% fields
\newcommand{\NN}{\mathbbmss{N}} 
\newcommand{\ZZ}{\mathbbmss{Z}} 
\newcommand{\QQ}{\mathbbmss{Q}} 
\newcommand{\RR}{\mathbbmss{R}} 
\newcommand{\CC}{\mathbbmss{C}} 
\newcommand{\KK}{\mathbbmss{K}} 
\newcommand{\FF}{\mathbbmss{F}} 
\newcommand{\HH}{\mathbbmss{H}} 

% symmetric group
\newcommand{\fS}{\mathfrak{S}}  

% calligraphic 
\newcommand{\aA}{\mathcal{A}} 
\newcommand{\bB}{\mathcal{B}}
\newcommand{\cC}{\mathcal{C}}
\newcommand{\dD}{\mathcal{D}}
\newcommand{\eE}{\mathcal{E}}
\newcommand{\fF}{\mathcal{F}}
\newcommand{\hH}{\mathcal{H}}
\newcommand{\iI}{\mathcal{I}}
\newcommand{\lL}{\mathcal{L}}
\newcommand{\oO}{\mathcal{O}}
\newcommand{\pP}{\mathcal{P}}
\newcommand{\qQ}{\mathcal{Q}}
\newcommand{\sS}{\mathcal{S}}
\newcommand{\mM}{\mathcal{M}}
\newcommand{\uU}{\mathcal{U}}

% bold
\newcommand{\bfa}{\mathbf{a}}
\newcommand{\bfe}{\mathbf{e}}
\newcommand{\bfF}{\pmb{F}}
\newcommand{\bfR}{\pmb{R}}
\newcommand{\bfv}{\mathbf{v}}
%\newcommand{\bfx}{\bm{x}}
%\newcommand{\bfx}{\mathbf{x}} 
\newcommand{\bfx}{\pmb{x}}
\newcommand{\bfX}{\pmb{X}}
\newcommand{\bfy}{\pmb{y}}
\newcommand{\bfz}{\pmb{z}}

% roman
\newcommand{\rmA}{\mathrm{A}}
\newcommand{\rmB}{\mathrm{B}}
\newcommand{\rmC}{\mathrm{C}}
\newcommand{\rmD}{\mathrm{D}} 
\newcommand{\rmF}{\mathrm{F}} 
\newcommand{\rmH}{\mathrm{H}} 
\newcommand{\rmh}{\mathrm{h}} 
\newcommand{\rmI}{\mathrm{I}} 
\newcommand{\rmK}{\mathrm{K}}
\newcommand{\rmM}{\mathrm{M}}
\newcommand{\rmP}{\mathrm{P}}  
\newcommand{\rmp}{\mathrm{p}}  
\newcommand{\rmQ}{\mathrm{Q}}  
\newcommand{\rmR}{\mathrm{R}}
\newcommand{\rmS}{\mathrm{S}}
\newcommand{\rmT}{\mathrm{T}}
\newcommand{\rmU}{\mathrm{U}}
\newcommand{\rmV}{\mathrm{V}}
\newcommand{\rmY}{\mathrm{Y}}
\newcommand{\rmZ}{\mathrm{Z}}
\newcommand{\rmz}{\mathrm{z}}

% arrows and symbols 
\renewcommand{\to}{\rightarrow}
\newcommand{\toto}{\longrightarrow}
\newcommand{\mapstoto}{\longmapsto}
\newcommand{\then}{\Rightarrow}
\newcommand{\IFF}{\Leftrightarrow}
\newcommand{\tl}{\tilde}
\newcommand{\wtl}{\widetilde}
\newcommand{\ol}{\overline}
\newcommand{\ul}{\underline}
\newcommand{\oldemptyset}{\emptyset}
\renewcommand{\emptyset}{\varnothing}
\DeclareMathSymbol{\Arg}{\mathbin}{AMSa}{"39} % for arguments 
\newcommand{\onto}{\ensuremath{\twoheadrightarrow}}
\newcommand{\tle}{\trianglelefteq}
\newcommand{\tge}{\trianglerighteq}

% custom ops
\newcommand{\rmKer}{\mathrm{Ker}}
\newcommand{\rmIm}{\mathrm{Im}}
\newcommand{\lcm}{\mathrm{lcm}}
\newcommand{\GL}{\mathrm{GL}}
\newcommand{\ev}{\mathrm{ev}}
\newcommand{\End}{\mathrm{End}}
\newcommand{\rmid}{\mathrm{id}}
\newcommand{\rmspan}{\mathrm{span}}
\newcommand{\Hom}{\mathrm{Hom}}
\newcommand{\Ann}{\mathrm{Ann}}
\newcommand{\Set}{\pmb{\mathrm{Set}}}
\newcommand{\Rmod}{\pmb{\mathrm{mod}}}
\newcommand{\rmrank}{\mathrm{rank}}
\newcommand{\mSpec}{\mathrm{mSpec}}
\newcommand{\Spec}{\mathrm{Spec}}

% absolute value symbol
\usepackage{mathtools} 
\DeclarePairedDelimiter\abs{\lvert}{\rvert}%
\DeclarePairedDelimiter\norm{\lVert}{\rVert}%
\makeatletter
\let\oldabs\abs
\def\abs{\@ifstar{\oldabs}{\oldabs*}}

% tensor symbol
\newcommand{\tensor}[1]{%
  \mathbin{\mathop{\otimes}\limits_{#1}}%
}

% permutation cycle notation
\ExplSyntaxOn
\NewDocumentCommand{\cycle}{ O{\;} m }
 {
  (
  \alec_cycle:nn { #1 } { #2 }
  )
 }

\seq_new:N \l_alec_cycle_seq
\cs_new_protected:Npn \alec_cycle:nn #1 #2
 {
  \seq_set_split:Nnn \l_alec_cycle_seq { , } { #2 }
  \seq_use:Nn \l_alec_cycle_seq { #1 }
 }
\ExplSyntaxOff

% setminus symbol
\newcommand{\mysetminusD}{\hbox{\tikz{\draw[line width=0.6pt,line cap=round] (3pt,0) -- (0,6pt);}}}
\newcommand{\mysetminusT}{\mysetminusD}
\newcommand{\mysetminusS}{\hbox{\tikz{\draw[line width=0.45pt,line cap=round] (2pt,0) -- (0,4pt);}}}
\newcommand{\mysetminusSS}{\hbox{\tikz{\draw[line width=0.4pt,line cap=round] (1.5pt,0) -- (0,3pt);}}}
\newcommand{\sm}{\mathbin{\mathchoice{\mysetminusD}{\mysetminusT}{\mysetminusS}{\mysetminusSS}}}

% miscellaneous commands
\newcommand{\defn}[1]{{\color{mylightblue}{#1}}}
\newcommand{\toDo}{{\bf\color{red} TODO}}
\newcommand{\toCite}{{\bf\color{green} CITE}}
\newcommand*{\vertbar}{\rule[-1ex]{0.5pt}{2.5ex}} % for matrices with column vectors
\newcommand*{\horzbar}{\rule[.5ex]{2.5ex}{0.5pt}} % for matrices with row vectors
\newcommand{\myblue}[1]{{\color{iceberg}{#1}}}
\newcommand{\myorange}[1]{{\color{burntorange}{#1}}}
\newcommand{\mygreen}[1]{{\color{applegreen}{#1}}}
\newcommand{\myred}[1]{{\color{darkcandyapplered}{#1}}}

% ferrer's diagram
\newcommand{\fdiagram}[1]{
    \begin{tikzpicture}[scale=.7]
        \fill foreach \Z [count=\Y] in {#1}
        {foreach \X in {1,...,\Z} 
        {(\X,-\Y) circle[radius=3pt]}};
    \end{tikzpicture}
}

%
\usepackage{pgfplots}
\pgfplotsset{compat=1.18}

%
\makeatletter
\def\mathcolor#1#{\@mathcolor{#1}}
\def\@mathcolor#1#2#3{%
  \protect\leavevmode
  \begingroup
    \color#1{#2}#3%
  \endgroup
}
\makeatother

% 
\newenvironment{nouppercase}{%
  \let\uppercase\relax%
  \renewcommand{\uppercasenonmath}[1]{}}{}

% titlepage
\title{226: Θεωρία Δακτυλίων και Modules}
\author[Β.~Δ. Μουστακας]{Βασίλης Διονύσης Μουστάκας \\ Πανεπιστήμιο Κρήτης}
\date{31 Μαρτίου 2026}
% \urladdr{\href{https://sites.google.com/view/vasmous}{https://sites.google.com/view/vasmous}}

\begin{document}

\begingroup
\def\uppercasenonmath#1{} % this disables uppercase title
\let\MakeUppercase\relax % this disables uppercase authors
\maketitle
\endgroup

\setcounter{section}{6}
\setcounter{theorem}{6}
\setcounter{equation}{4}
\begin{center}
    \textbf{6. Τανυστικό γινόμενο
} (Συνέχεια)
\end{center}

Υπενθυμίζουμε πως σε ότι ακολουθεί, $R$ είναι ένας μεταθετικός δακτύλιος.

Η επόμενη πρόταση μας επιτρέπει να \textquote{αναγνωρίσουμε} τανυστικά γινόμενα ελεύθερων προτύπων (και ειδικότερα διανυσματικών χώρων) και ολοκληρώνει τον στόχο μας να κατασ\-κευάσουμε ένα \textquote{γινόμενο} δύο προτύπων.

\begin{proposition}
    \label{prop:tensorProduct_free}
    Έστω $M$ και $N$ δύο ελεύθερα $R$-πρότυπα με βάσεις $\{a_\lambda : \lambda \in \Lambda\}$ και $\{b_\kappa : \kappa \in K\}$, αντίστοιχα. Το σύνολο 
    \[
    \{a_\lambda \otimes b_\kappa : \lambda \in \Lambda, \ \kappa \in K\}
    \]
    αποτελεί βάση του $M \otimes_R N$. Ειδικότερα,
    \[
    \rmrank_R(M\otimes_RN) = \rmrank_R(M)\rmrank_R(N).
    \]
\end{proposition}

\begin{proof}[Απόδειξη]
    Αρχικά, παρατηρούμε ότι το σύνολο 
    \[
    \{a_\lambda \otimes b_\kappa : \lambda \in \Lambda, \ \kappa \in K\}
    \]
    παράγει το $M \otimes_R N$. Πράγματι, αν $a \in M$ και $b \in N$, τότε 
    \begin{equation}
        \label{eq:tensor_help}
        a = \sum_{\lambda \in \Lambda} r_\lambda a_\lambda \quad \text{και} \quad 
        b = \sum_{\kappa \in K} s_\kappa b_\kappa
    \end{equation}
    για κάποια $r_\lambda = s_\kappa = 0_R$ εκτός από το πολύ ένα πεπερασμένο πλήθος. Συνεπώς, 
    \[
    a \otimes b = \left(
        \sum_{\lambda \in \Lambda} r_\lambda a_\lambda
    \right) \otimes 
    \left(
        \sum_{\kappa \in K} s_\kappa b_\kappa
    \right) = \sum_{\substack{\lambda \in \Lambda \\ \kappa \in K}} r_\lambda s_\kappa \, a_\lambda \otimes b_\kappa,
    \]
    όπου $r_\lambda s_\kappa = 0_R$ εκτός από το πολύ ένα πεπερασμένο πλήθος.

    Μένει να δείξουμε ότι είναι $R$-γραμμικώς ανεξάρτητο. Για τον λόγο αυτό, υποθέτουμε ότι 
    \begin{equation}
        \label{eq:tensor_helpHelp}
        \sum_{\substack{\lambda \in \Lambda \\ \kappa \in K}} r_{\lambda\kappa} \, a_\lambda \otimes b_\kappa = 0_{M\otimes_R{N}},
    \end{equation}
    για κάποια $r_{\lambda\kappa} = 0_R$ εκτός από το πολύ ένα πεπερασμένο πλήθος. Θα δείξουμε ότι $r_{\lambda\kappa} = 0_R$ για κάθε $\lambda \in \Lambda$ και $\kappa \in K$. Για κάθε $\lambda_0 \in \Lambda$ και $\kappa_0 \in K$, θεωρούμε την απεικόνιση 
    \begin{align*}
        \phi_0 : M \times N &\to R \\
        (a, b) &\mapsto r_{\lambda_0} s_{\kappa_0},
    \end{align*}
    όπου τα $a$ και $b$ αναπτύσονται στις αντίστοιχες βάσεις όπως στην \eqref{eq:tensor_help}. Η $\phi$ είναι $R$-διγραμμική και γι αυτό επάγει μία $R$-γραμμική απεικόνιση $\phi_0 : M \otimes_R N \to R$ τέτοια ώστε 
    \[
    \phi_0(a\otimes{b}) = r_{\lambda_0} s_{\kappa_0}.
    \]
    Ειδικότερα, 
    \[
    \phi_0(a_\lambda \otimes b_\kappa) = 
    \begin{cases}
        1_R, &\ \text{αν $(\lambda,\kappa) = (\lambda_0,\kappa_0)$} \\
        0_R, &\ \text{διαφορετικά}.
    \end{cases}
    \]
    Εφαρμόζοντας την $\phi_0$ στην σχέση \eqref{eq:tensor_helpHelp} προκύπτει ότι $r_{\lambda_0\kappa_0} = 0_R$ και το ζητούμενο έπεται.
\end{proof}

\begin{example_6.6}{\rm(Συνέχεια)}
    \leavevmode
    \begin{itemize}
        \item[(η)] Έχουμε 
        \[
        \rmM_{nm}(R) \otimes_R \rmM_{k\ell}(R) \cong \rmM_{(nk)(m\ell)}(R).
        \]
        Όμοια με το Παράδειγμα 5.2 (δ), το σύνολο $\{E_{i,j}^{n,m} : 1 \leq i \leq n, \ 1 \le j \le m\}$ (αντ. $\{E_{i,j}^{k,\ell} : 1 \leq i \leq k, \ 1 \le j \le \ell\}$) αποτελεί βάση του $\rmM_{nm}(R)$ (αντ. $\rmM_{k\ell}(R)$). Από την Πρόταση \ref{prop:tensorProduct_free} έπεται ότι το σύνολο 
        \[
        \{
            E_{i,j}^{n,m} \otimes E_{s,t}^{k,\ell} : 1 \leq i \leq n, \ 1 \le j \le m, \ 1 \leq s \leq k, \ 1 \le t \le \ell
        \}
        \]
        αποτελεί βάση του $\rmM_{nm}(R) \otimes_R \rmM_{k\ell}(R)$. Η απεικόνιση 
        \begin{align*}
            \rmM_{nm}(R) \otimes_R \rmM_{k\ell}(R) &\to \rmM_{(nk)(m\ell)}(R) \\ 
            E_{i,j}^{nm} \otimes E_{s,t}^{k\ell} &\mapsto 
            E_{(i-1)k +s, (j-1)\ell+t}^{nk, m\ell}
        \end{align*}
        είναι $R$-ισμομορφισμός (γιατί;).

        Για παράδειγμα, για $n=2, m=3, k=\ell=2$ έχουμε 
        \[
        \begin{pmatrix}
        0 & 0 & 0 \\
        0 & 1 & 0
        \end{pmatrix}
         \otimes \begin{pmatrix}
            0 & 1 \\
            0 & 0
         \end{pmatrix} \mapsto
        \left(
        \begin{array}{cc|cc|cc}
        0 & 0 & 0 & 0 & 0 & 0 \\
        0 & 0 & 0 & 0 & 0 & 0 \\
        \hline
        0 & 0 & 0 & 1 & 0 & 0 \\
        0 & 0 & 0 & 0 & 0 & 0
        \end{array}
        \right)
        \]

        Γενικότερα, αν $A = (a_{ij})_{\substack{1 \le i \le n\\ 1 \le j \le m}}$, τότε
        \[
        A \otimes B \mapsto
        \begin{pmatrix}
            a_{11}B & a_{12}B & \cdots & a_{1m}B \\
            a_{21}B & a_{22}B & \cdots & a_{2m}B \\
            \vdots & \vdots &        & \vdots \\
            a_{n1}B & a_{n2}B & \cdots & a_{nm}B
        \end{pmatrix} 
        \]
        όπου στο δεξί μέλος έχουμε τον μπλόκ-πίνακα του οποίου το $(i,j)$-μπλοκ είναι ο πίνακας $a_{ij}B$. Ο πίνακας αυτός ονομάζεται \emph{γινόμενο Kronecker} των $A$ και $B$.
        \item[(θ)] Έχουμε 
        \[
        R[x] \otimes_R R[y] \cong R[x,y].
        \]
        Πράγματι, στο Παράδειγμα 5.2 (στ) είδαμε ότι το σύνολο $\{1, x, x^2, \dots\}$ αποτελεί βάση για το $R[x]$ ως $R$-πρότυπο. Συνεπώς, από την Πρόταση \ref{prop:tensorProduct_free} έπεται ότι το σύνολο 
        \[
        \{x^i \otimes y^j : i, j \ge 0\}
        \]
        αποτελεί βάση για το $R[x] \otimes_R R[y]$. Η απεικόνιση
        \begin{align*}
            R[x] \otimes_R R[y] &\to R[x,y] \\ 
            x^i \otimes y^j &\mapsto x^iy^j
        \end{align*}
        είναι $R$-ισομορφισμός (γιατί;). 
    \end{itemize}
\end{example_6.6}

Ο ισομορφισμός του Παραδείγματος 6.6 (θ) ισχύει στο επίπεδο των $R$-αλγεβρών.
% \ref{ex:tensorProduct}
\begin{definition}
    \label{def:algebras}
    Ένας δακτύλιος $A$ εφοδιασμένος με έναν ομομορφισμό δακτυλίων $u : R \to A$ τέτοιο ώστε 
    \[     
    u(R) \subseteq \rmZ(A) \coloneqq \{a \in A : ab=ba, \ \text{για κάθε $b \in A$}\}
    \]
    ονομάζεται $R$-άλγεβρα.
\end{definition}

Κάθε $R$-άλγεβρα έχει δομή $R$-προτύπου με δράση 
\[
r\cdot{a} \coloneqq u(r)a
\]
για κάθε $r \in R$ και $a \in A$. Η συνθήκη $u(R) \subseteq \rmZ(A)$ μας πληροφορεί ότι 
\begin{equation}  
    \label{eq:algebra}
    r\cdot(ab) = (r\cdot{a})b = a(r\cdot{b}).
\end{equation}
Ισχύει και το αντίστροφο (γιατί;). Με άλλα λόγια, οι $R$-άλγεβρες είναι δκατύλιοι που έχουν δομή $R$-προτύπου έτσι ώστε ο πολλαπλασιασμός και η δράση να ικανοποιούν την συνθήκη \eqref{eq:algebra}.

\begin{example}
    \label{ex:algebra}
    \leavevmode
    \begin{itemize}
        \item[(α)] Κάθε δακτύλιος είναι $\ZZ$-άλγεβρα.
        \item[(β)] Ο $\rmM_n(R)$ είναι $R$-άλγεβρα με $u(1_R) = I_n$.
        \item[(γ)] Ο $R[x_1, x_2, \dots, x_n]$ είναι $R$-άλγεβρα με $u(1_R) = 1_R$.  Γενικότερα, κάθε μεταθετικός δακτύλιος $R$ είναι $S$-άλγεβρα, για κάθε υποδακτύλιο $S \subseteq R$.  
        \item[(δ)] Κάθε δακτύλιος πηλίκο $R[x_1, x_2, \dots, x_n]/I$ είναι $R$-άλγεβρα.
        \item[(ε)] Έστω $M$ ένα $R$-πρότυπο. Ο $\End_R(M)$ είναι $R$-άλγεβρα με $u(1_R) = \rmid_M$.
    \end{itemize}
\end{example}

Όπως δείχνει η Ταυτότητα (6.4),
% \eqref{eq:tensor_associativity} 
μπορούμε να ορίσουμε το τανυστικό γινόμενο 
\[
M_1 \otimes M_2 \otimes \cdots \otimes M_k
\]  
ενός πεπερασμένου αριθμού προτύπων $M_1, M_2, \dots, M_k$. Αυτό ικανοποιεί μία καθολική ιδιό\-τητα παρόμοια με αυτή του Θεωρήματος 6.3.
% \ref{thm:UP_tensorProduct}

Πιο συγκεκριμένα, για κάθε $R$-πρότυπο $P$ και κάθε \emph{$R$-πολυγραμμική} ($R$-multilinear) απει\-κόνιση $\phi : M_1 \times M_2 \times \cdots \times M_k \to P$, δηλαδή η $\phi$ είναι $R$-γραμμική στην $i$-οστή συντεταγμένη με τις υπόλοιπες σταθεροποιημένες, για κάθε $1 \le i \le k$, υπάρχει μοναδική $R$-γραμμική απεικόνιση $\tilde{\phi} : M_1 \otimes M_2 \otimes \cdots \otimes M_k \to P$ τέτοια ώστε το διάγραμμα 
\[
    \begin{tikzcd}
    M_1 \times M_2 \times \cdots \times M_k \arrow[r, "\otimes^k"] \arrow[dr, "\phi"'] & M_1 \otimes M_2 \otimes \cdots \otimes M_k \arrow[d, "\tilde{\phi}"] \\
                                        & P
    \end{tikzcd}
\]
να είναι μεταθετικό, ή ισοδύναμα, 
\[
\tilde{\phi}(a_1 \otimes a_2 \otimes \cdots \otimes a_k) = \phi(a_1,a_2,\dots, a_k),
\]
για κάθε $a_i \in M_i$, για κάθε $1 \le i \le k$.

\begin{proposition}
    \label{prop:tensorProduct_algebras}
    Έστω $A$ και $B$ δύο $R$-άλγεβρες. Το $A \otimes_R B$ με πολλαπλασιασμό
    \[
    (a \otimes b) (a' \otimes b') \coloneqq (aa') \otimes (bb')
    \]
    είναι $R$-άλγεβρα.
\end{proposition}

\begin{proof}[Απόδειξη]
    Η επαλήθευση της συνθήκης \eqref{eq:algebra} αφήνεται ως άσκηση. Μένει να δείξουμε ότι ο πολλαπλασιασμός είναι καλώς ορισμένος. Για τον λόγο αυτό, θεωρούμε την απεικόνιση 
    \begin{align*}
        \phi : A \times B \times A \times B &\to A \otimes B \\
        (a,b,a',b') &\mapsto (aa') \otimes (bb').
    \end{align*}
    Αυτή είναι $R$-πολυγραμμική και από τη συζήτηση πριν την Πρόταση \ref{prop:tensorProduct_algebras} επάγει μία $R$-γραμμική απεικόνιση 
    \begin{align*}
       \tilde{\phi} : A \otimes B \otimes A \otimes B &\to A \otimes B \\ 
        a\otimes b \otimes a' \otimes b' &\mapsto (aa') \otimes (bb').
    \end{align*}
    Σκεπτόμενοι το τανυστικό γινόμενο $A \otimes B \otimes A \otimes B$ ως $(A \otimes B) \otimes_R (A \otimes B)$ (λόγω προσεταιριστικότητας), η $\tilde{\phi}$ αντιστοιχεί σε μία $R$-διγραμμική απεικόνιση 
    \begin{align*}
        (A \otimes B) \times (A \otimes B) &\to A \otimes B \\
        \left(
            (a \otimes b), (a' \otimes b')
        \right) &\mapsto \tilde{\phi}(a\otimes b \otimes a' \otimes b') = (aa') \otimes (bb').
    \end{align*}
\end{proof}

\begin{example}
    \label{ex:tensorProduct_algebras}
    Για κάθε $R$-άλγεβρα $A$, 
    \[
    A \otimes_R R[x] \cong A[x],
    \]
    ως $R$-άλγεβρες. Ειδικότερα, 
    \[
    R[x] \otimes_R R[y] \cong \left(
        R[x]
    \right)[y] = R[x,y].
    \]
    Πράγματι, υπενθυμίζουμε ότι το σύνολο $\{x^i : i \ge 0\}$ αποτελεί βάση του $R[x]$ ως $R$-πρότυπο. Με αυτό κατά νου, θεωρούμε την απεικόνιση 
    \begin{align*}
        \phi : A \times_R R[x] &\to A[x] \\
        (a, x^i) &\mapsto ax^i.
    \end{align*}
    Η $\phi$ είναι $R$-διγραμμική (γιατί;) και γι αυτό επάγει μία $R$-γραμμική απεικόνιση $\tilde{\phi} : A \otimes_R R[x] \to A[x]$ τέτοια ώστε $\tilde{\phi}(a \otimes x^i) = ax^i$. Επιπλέον,
    \begin{align*}
    \tilde{\phi}\left(
        \left(
            a \otimes x^i
        \right)\left(
            a' \otimes x^j
        \right)
    \right) &= \tilde{\phi}\left(
        (aa') \otimes x^{i+j}
    \right) \\
    &= aa'x^{i+j} \\
    &= (ax^i)(a'x^j) \\
    &= \tilde{\phi}(a\otimes x^i)\tilde{\phi}(a' \otimes x^j),
    \end{align*}
    για κάθε $a, a' \in A$ και $i, j \ge 0$. Συνεπώς, η $\tilde{\phi}$ είναι και ομομορφισμός δακτυλίων και κατά συνέπεια ομομορφισμό $R$-αλγεβρών. Για την αντίστροφή της, θεωρούμε την απεικόνιση 
    \begin{align*}
        \psi : A[x] &\to A \otimes_R R[x] \\
            ax^i &\mapsto a \otimes x^i
    \end{align*}
    Αυτή είναι ομομορφισμός $R$-αλγεβρών, δηλαδή $R$-γραμμική και ομομορφισμός δακτυλίων (γιατί;). Επιπλέον,
    \[
        \tilde{\phi}\circ\psi = \rmid_{A[x]} \quad \text{και} \quad 
        \psi\circ\tilde{\phi} = \rmid_{A \otimes_R R[x]}
    \]
    (γιατί;) και το ζητούμενο έπεται.
\end{example}


% \begin{thebibliography}{99}
%     \bibitem{Alu09}
%     P.~Aluffi,
%     \textit{Algebra, Chapter 0},
%     Grad. Stud. in Math. {\bf104},
%     Amer. Math. Soc., Providence, RI, 2009.
%     \bibitem{DF04}
%     D. S. Dummit και R. M. Foote, 
%     \textit{Abstract algebra},
%     third edition, Wiley, 2004. 
% \end{thebibliography}
\end{document}